\pagestyle{mystyle}
\chapter{Valutazione Sperimentale e Analisi dei Risultati}
\label{chap:evaluation}

Il capitolo definisce la valutazione sperimentale del sistema dopo le scelte
metodologiche descritte nel Capitolo~\ref{chap:ottimizzazione}. Poiché i
benchmark completi non sono ancora stati eseguiti, il testo distingue tra
protocollo di valutazione, metriche da raccogliere e struttura attesa
dell'analisi. Questa distinzione è necessaria per evitare di presentare come
risultati misurati dati che appartengono ancora alla fase di pianificazione
sperimentale.

La domanda di ricerca è la seguente: in che misura la quantizzazione di Gemma 3
4B Instruct tramite MLX modifica il compromesso tra memoria, latenza, throughput
e qualità generativa su Apple Silicon consumer, e quali configurazioni risultano
sostenibili per una pipeline RAG enterprise on-premise? La risposta richiede due
livelli di analisi. Il primo riguarda il modello isolato, misurato in termini di
prestazioni e qualità generale. Il secondo riguarda il sistema RAG completo,
dove la qualità non dipende solo dal modello generativo, ma anche dal retrieval,
dalla costruzione del prompt e dalla capacità di restare aderente alle fonti.

Il confronto principale è quindi interno a MLX: la variabile sperimentale è il
livello di quantizzazione del modello, mentre runtime, hardware, prompt e
parametri di generazione restano controllati. GGUF, llama.cpp e Ollama non
vengono rimossi dalla valutazione, ma assumono un ruolo diverso: costituiscono
un riferimento applicativo secondario, utile per collegare l'esperimento al
prototipo e alla riproducibilità pratica, non la base del confronto scientifico
su Apple Silicon.

\section{Ambiente sperimentale}
\label{sec:evaluation_environment}

L'ambiente di riferimento è un MacBook Pro con chip M1 Pro, 16 GB di unified
memory, macOS e backend Metal. Il vincolo hardware è parte dell'esperimento:
non si valuta una workstation specializzata, ma una macchina consumer che può
rappresentare un limite realistico per prototipi aziendali on-premise. Questa
scelta rende il confronto più vicino al caso d'uso della tesi, anche se limita
la trasferibilità dei risultati verso sistemi con più memoria o GPU dedicate.

Le configurazioni principali sono la baseline MLX in precisione non quantizzata
e più varianti quantizzate generate con \texttt{mlx-lm}. La granularità esatta
dei livelli dipende dal supporto disponibile nel runtime al momento
dell'esecuzione, ma il protocollo prevede almeno un confronto tra baseline,
quantizzazione moderata e quantizzazione aggressiva. Se tecnicamente
disponibili, verranno incluse configurazioni a 8, 6 e 4 bit. La configurazione
GGUF/Ollama, eventualmente eseguita su una o due varianti equivalenti, viene
riportata separatamente e non viene mescolata con il confronto principale MLX.

I parametri di generazione restano costanti tra le configurazioni: contesto pari
a 4096 token, massimo 256 token generati, temperatura 0, \texttt{top\_p} pari a
1 e seed fisso quando supportato dal runtime. Anche corpus, embedding,
retrieval, valore di \texttt{top\_k} e template di prompt restano invariati nei
test RAG. In questo modo il cambiamento osservato può essere ricondotto alla
quantizzazione e non a modifiche concorrenti della pipeline.

\begin{table}[ht]
  \centering
  \caption{Ruolo delle configurazioni nella valutazione sperimentale.}
  \label{tab:evaluation_configurations}
  \small
  \begin{tabular}{p{3cm}p{3cm}p{6.5cm}}
    \toprule
    \textbf{Configurazione} & \textbf{Ruolo} & \textbf{Uso nell'analisi} \\
    \midrule
    MLX FP16/BF16 & Baseline principale
      & Riferimento per qualità, memoria, latenza e throughput. \\
    MLX quantizzata moderata & Confronto principale
      & Misura del compromesso tra compressione e qualità. \\
    MLX quantizzata aggressiva & Confronto principale
      & Valutazione del limite inferiore di memoria e del degrado qualitativo. \\
    GGUF/Ollama & Baseline applicativa secondaria
      & Collegamento con il prototipo e con strumenti diffusi di deployment locale. \\
    \bottomrule
  \end{tabular}
\end{table}

\section{Metriche prestazionali}
\label{sec:evaluation_performance_metrics}

Il benchmark prestazionale raccoglie metriche che descrivono sia il costo di
deployment sia il comportamento interattivo. Per ogni variante vengono
registrati dimensione su disco, rapporto di compressione rispetto alla baseline,
riduzione percentuale della dimensione, tempo di caricamento, memoria residente,
time to first token, tempo totale di generazione e token al secondo. Questi
valori non devono essere letti separatamente: un modello più piccolo può ridurre
la memoria occupata ma non necessariamente migliorare la latenza, perché su
Apple Silicon il costo di dequantizzazione, l'accesso alla unified memory e
l'efficienza dei kernel Metal possono pesare in modo rilevante
~\cite{benazir2025profiling,rajesh2025productionlocal}.

Per rendere il confronto più utile in ottica enterprise vengono calcolate anche
metriche normalizzate, come token al secondo per GiB e latenza per dimensione
del modello. Questo tipo di misura aiuta a individuare configurazioni efficienti
rispetto alla memoria disponibile, non solo rispetto alla velocità assoluta. In
un sistema on-premise con risorse limitate, infatti, la configurazione migliore
non coincide necessariamente con quella più veloce o più compressa, ma con
quella che offre stabilità, qualità e consumo di memoria compatibili con l'uso
applicativo.

\section{Metriche di qualità}
\label{sec:evaluation_quality_metrics}

La qualità viene misurata su tre livelli. Il primo è la perplexity su un testo
di riferimento, utile per osservare il degrado statistico introdotto dalla
quantizzazione. Il secondo usa subset ridotti di MMLU e GSM8K, per ottenere un
segnale, anche limitato, su conoscenza generale e ragionamento matematico. Il
terzo livello è specifico per la tesi: un benchmark RAG costruito su un corpus
enterprise simulato, con domande che richiedono risposte fondate sui documenti
recuperati.

Nel benchmark RAG la valutazione considera groundedness, completezza,
allucinazioni e latenza end-to-end. La groundedness misura quanto la risposta
sia supportata dal contesto recuperato; la completezza misura se copre le
informazioni richieste; il conteggio delle allucinazioni segnala affermazioni
non giustificate dalle fonti. Per collegare qualità e costo di deployment,
queste metriche possono essere normalizzate rispetto alla dimensione del
modello, ad esempio groundedness per GiB o completezza per GiB.

\begin{table}[ht]
  \centering
  \caption{Metriche previste per la valutazione RAG.}
  \label{tab:rag_evaluation_metrics}
  \small
  \begin{tabular}{p{3.5cm}p{2.2cm}p{7cm}}
    \toprule
    \textbf{Metrica} & \textbf{Scala} & \textbf{Significato} \\
    \midrule
    Groundedness & 1--5 & Aderenza della risposta al contesto recuperato. \\
    Completezza & 1--5 & Copertura delle informazioni richieste dalla domanda. \\
    Allucinazioni & 0--3 & Presenza di affermazioni non supportate dalle fonti. \\
    Latenza & secondi & Tempo end-to-end della risposta RAG. \\
    Qualità per GiB & derivata & Qualità normalizzata rispetto alla dimensione del modello. \\
    \bottomrule
  \end{tabular}
\end{table}

\section{Valutazione della generazione di bozze}
\label{sec:evaluation_draft_generation}

La generazione di bozze di Gilda richiede una valutazione separata, perché il
risultato non è una risposta libera ma un oggetto strutturato. Ogni output deve
essere verificato rispetto allo schema \texttt{GuildDraft}: presenza dei campi
obbligatori, coerenza dei tipi, leggibilità dei contenuti e compatibilità con il
flusso applicativo. Un output discorsivo o un JSON incompleto non viene contato
come successo, anche se il testo contiene informazioni plausibili.

Il tasso di errore deve essere riportato come numero di bozze non valide sul
totale dei test eseguiti. Questa metrica è importante perché il sistema non usa
un fallback artificiale: se la risposta non è valida, l'utente viene informato
dell'errore e invitato a riprovare. La valutazione deve quindi misurare non solo
la qualità media delle bozze riuscite, ma anche la probabilità che il modello
produca un output utilizzabile al primo tentativo.

\section{Struttura dell'analisi dei risultati}
\label{sec:evaluation_result_structure}

Quando i benchmark saranno eseguiti, l'analisi dovrà confrontare ogni variante
MLX contro la baseline non quantizzata nello stesso runtime. Per ciascuna
configurazione verranno riportati delta di memoria, delta di velocità, delta di
qualità e metriche normalizzate. Una lettura corretta deve evitare conclusioni
monotone: la quantizzazione più aggressiva potrebbe essere la variante più
piccola, ma non per questo la più utile; una quantizzazione leggera potrebbe
mantenere qualità vicina alla baseline, ma senza offrire il miglior rapporto tra
qualità e memoria; una quantizzazione intermedia potrebbe risultare il
compromesso migliore, ma solo i dati locali possono confermarlo.

Il confronto GGUF/Ollama, se eseguito, dovrà essere discusso in una sottosezione
separata. Il suo scopo non è stabilire in modo generale quale runtime sia
migliore, ma verificare se le conclusioni ottenute nel setup MLX rimangono
coerenti con l'integrazione applicativa usata dal prototipo. Questa separazione
evita di attribuire alla sola quantizzazione differenze prodotte dal runtime,
dal formato del modello o dai rispettivi backend di inferenza.

La parte finale dovrà discutere i casi di fallimento: documenti non recuperati,
query ambigue, risposte corrette ma incomplete, risposte non sufficientemente
ancorate alle fonti, JSON non validi nella generazione delle bozze e limiti dei
modelli locali piccoli. Questo passaggio collega la valutazione tecnica al caso
d'uso enterprise: un modello sostenibile non è solo un modello che entra in
memoria, ma un modello che mantiene un comportamento affidabile nel flusso RAG
effettivamente usato dall'applicazione.
